\chapter{Conclusion}
\label{ch:conclusion}

Ce projet est parti d'une question simple mais exigeante : sur un réseau routier
réel, à réponse optimale identique, jusqu'où de meilleures heuristiques et de
meilleures structures de données peuvent-elles réduire le travail d'un calcul de
plus court chemin ? Pour y répondre sérieusement, nous avons d'abord modélisé une
carte OpenStreetMap de Paris en un graphe orienté pondéré, avec les choix
d'ingénierie nécessaires pour opérer à l'échelle, puis nous avons construit une
échelle complète d'algorithmes autour de ce même graphe.

\medskip

Cette échelle constitue le cœur du travail. Elle va des parcours non informés
(BFS, DFS) et de la recherche gloutonne aux méthodes optimales fondées sur une
heuristique, A* puis ALT, avant de dépasser les fondamentaux avec la recherche
bidirectionnelle, les files à seaux, les Arc Flags et les Contraction
Hierarchies. Chaque algorithme a été non seulement implémenté, mais justifié
mathématiquement et vérifié comme rendant exactement le plus court chemin. Autour
de ces algorithmes, nous avons mis en place un véritable dispositif expérimental :
un banc d'essai agrégé, une étude du placement des landmarks, une étude de passage
à l'échelle, une analyse statistique de la difficulté des requêtes, ainsi qu'un
trajet réel à plusieurs étapes accompagné de son itinéraire détaillé, d'une carte
interactive et d'une animation de l'exploration.

\medskip

Les résultats confirment la théorie de bout en bout. Ajouter de l'information, de
la structure et du pré-calcul réduit l'exploration de plus d'un ordre de grandeur
sans jamais sacrifier l'optimalité, tandis que les algorithmes non informés et
glouton, placés en regard, rappellent pourquoi la pondération et l'optimalité
comptent. Surtout, ces gains ne sont pas figés : ils s'amplifient à mesure que le
réseau grandit. C'est là l'enseignement central du projet : la difficulté est
l'échelle, et la réponse est d'investir une fois dans un pré-calcul pour payer
beaucoup moins à chaque requête, un compromis d'autant plus rentable que la carte
est grande, exactement le principe sur lequel reposent les moteurs de routage
modernes.

\medskip

Plusieurs pistes prolongeraient naturellement ce travail : une implémentation
compilée pour révéler pleinement l'avantage des méthodes hiérarchiques, la prise
en compte de coûts dépendant du temps pour modéliser le trafic, ou encore le
passage à un réseau de taille continentale. Au-delà des algorithmes particuliers,
le projet illustre une idée générale et durable : face à un très grand espace
d'états, ce n'est pas la force brute qui l'emporte, mais l'information et le
pré-calcul judicieusement investis.
